% Generated draft: requires mathematical and source review.
\documentclass[11pt]{article}
\usepackage[margin=25mm]{geometry}
\usepackage{fontspec}
\setmainfont{Latin Modern Roman}
\usepackage{amsmath,amssymb,mathtools,unicode-math}
\setmathfont{Latin Modern Math}
\usepackage{enumitem,needspace,xurl,hyperref}
\setlength{\parindent}{0pt}
\setlength{\parskip}{5pt}
\setlength{\emergencystretch}{3em}
\newcommand{\N}{\mathbb{N}}
\newcommand{\Z}{\mathbb{Z}}
\newcommand{\Q}{\mathbb{Q}}
\newcommand{\R}{\mathbb{R}}
\newcommand{\C}{\mathbb{C}}
\newcommand{\F}{\mathbb{F}}
\newcommand{\A}{\mathbb{A}}
\newcommand{\PP}{\mathbb{P}}
\newcommand{\E}{\mathbb{E}}
\newcommand{\eps}{\varepsilon}
\newcommand{\dd}{\,\mathrm{d}}
\begin{document}
\section*{UnsolvedMath ID: 10400218}
\textbf{Source problem code:} catalogue entry
\subsection{Definitions and mathematical statement}
J. Roberts asks the foundational question: ``What are quantum groups?'' The excerpt gives no specified classification target, construction, universal property or theorem to establish.

Standard examples include Drinfeld--Jimbo Hopf-algebra deformations of universal enveloping algebras. In a formal deformation convention, one has a Hopf algebra \(U_\hbar(\mathfrak g)\) with classical limit \(U_\hbar(\mathfrak g)/\hbar U_\hbar(\mathfrak g)\cong U(\mathfrak g)\). Other usages concern deformations of function algebras or related structures. These examples explain the supplied summary's context but do not select a unique definition answering Roberts's question.

The source does not state whether it seeks an intrinsic characterization, geometric or topological meaning, classification, or a unifying framework for those usages. It must not be silently changed into the known construction of one quantum enveloping algebra, the assertion that every Hopf algebra is a quantum group, or a demand to classify all Hopf algebras. The qualitative question is retained, with its intended conceptual success criterion unresolved.
\subsection{Short English statement}
Roberts asks for conceptual understanding of quantum groups. The source provides no precise theorem or criterion, and standard deformation examples do not determine the full intended question.
\subsection{Sources}
\begin{itemize}
\item Local UnsolvedMath registry, permanent ID 10400218, source code catalogue entry; statement attribution J. Roberts; proposer field and proposal year absent; update 2026-08-06.
\item Ohtsuki, \emph{Problems on invariants of knots and 3-manifolds} (2002), Problem 12.13, PDF page 165: \url{https://msp.org/gtm/2002/04/gtm-2002-04-024s.pdf}. Registry access date 2026-07-29; PDF font-structure extraction. Provenance distinguishes 231 numbered bold headings. Present primary-literature status and redistribution terms require review.
\item The supplied summary credits Drinfeld and Jimbo for enveloping-algebra deformations and describes that as partial conceptual progress, without exact references or a formal resolution criterion. The registry's open-triage label is not independently assessed here.
\end{itemize}
\end{document}
